\documentclass[12pt]{article}
\usepackage[margin=1in]{geometry}
\usepackage{graphicx}
\usepackage{booktabs}
\usepackage{amsmath}
\usepackage{natbib}
\usepackage{setspace}
\onehalfspacing

\title{Replicating Tables 2 and 3 from \\ Luo and Spindler (2017)}
\author{Juan Villa}
\date{\today}

\begin{document}
\maketitle

% =============================================================================
\section{The Original Paper}
% =============================================================================

This replication targets \citet{luo2017l2boosting}, ``L2-Boosting for Economic
Applications,'' published in the \textit{American Economic Review: Papers \&
Proceedings} (2017, 107(5): 270--273). The paper asks whether L2-Boosting can serve as a competitive alternative to
Lasso for causal inference in high-dimensional settings, where the number of
potential covariates or instruments is large relative to the number of
observations.

The empirical strategy builds on the orthogonal moment conditions framework
of \citet{chernozhukov2016double}. Valid post-selection inference requires
that the variable selection step achieve a sufficient rate of convergence.
Luo and Spindler demonstrate, building on \citet{luo2016highdim}, that
post-L2-Boosting and orthogonal L2-Boosting satisfy these theoretical
conditions and achieve the same convergence rate as Lasso. They illustrate
the methods through two applications: IV estimation with many instruments
using the EminentDomain dataset, and inference on a treatment effect with
many controls using macroeconomic growth data.

% =============================================================================
\section{The Result of Interest}
% =============================================================================

This replication targets two specific results from the paper.

\textbf{Table 2} reports the estimated effect of federal appellate court
property rights decisions on log GDP per capita across four estimators.
The paper reports the following results:

\begin{center}
\begin{tabular}{lcccc}
\hline
 & post-Lasso & BA & post-BA & oBA \\
\hline
$\hat{\beta}$ & 0.005 & 0.005 & 0.004 & 0.008 \\
s.e.          & 0.012 & 0.007 & 0.006 & 0.006 \\
\hline
\end{tabular}
\end{center}

\textbf{Table 3} reports Monte Carlo simulation results evaluating estimator
performance under a controlled data-generating process with known true effect
$\alpha = 1$. The paper reports:

\begin{center}
\begin{tabular}{lcccc}
\hline
 & post-Lasso & BA & post-BA & oBA \\
\hline
bias & 0.082 & 0.121 & 0.136 & 0.121 \\
RP   & 0.002 & 0.042 & 0.054 & 0.042 \\
\hline
\end{tabular}
\end{center}

These two tables are the natural replication targets because together they
tell the paper's complete story: Table 3 evaluates the methods under
controlled conditions where the truth is known, and Table 2 shows their
performance on a real empirical problem. Table 2 is also the headline
empirical result of the paper and the most-cited output.

% =============================================================================
\section{Data and Methods}
% =============================================================================

\subsection*{Data}

Table 2 uses the \texttt{EminentDomain} dataset from the \texttt{hdm} R
package \citep{hdm2016}, originally constructed by \citet{belloni2012sparse}.
The dataset contains $n = 312$ observations with $p = 140$ potential
instrumental variables capturing federal appellate court judge characteristics
and court composition. The outcome is log GDP per capita. The treatment is the share of pro-plaintiff appellate court
decisions, measuring the strength of property rights protection.

Table 3 uses no real data. Observations are generated programmatically by
the DGP defined in \texttt{code/DGP.R}, taken directly from the authors'
replication package. The simulation draws $R = 500$ replications with
$n = 100$ observations and $p = 100$ potential instruments per draw.
The true treatment effect is $\alpha = 1$, the first stage is sparse with
$s = 5$ nonzero coefficients, and instruments follow an AR(1)-type covariance
structure with parameter 0.5. The target first-stage F-statistic is 30 and
errors are independent ($C_{ev} = 0$).

\subsection*{What \texttt{preprocess.R} Does}

The preprocessing script loads \texttt{EminentDomain} directly from the
\texttt{hdm} package. It extracts the outcome (\texttt{y}), treatment
(\texttt{d}), instruments (\texttt{z}), and controls (\texttt{x}), dropping
column 50 of the control matrix following the original replication code.
It then standardizes all variables: the outcome is demeaned only, while the
treatment, instruments, and controls are demeaned and scaled to unit variance.
The processed dataset is saved to \texttt{temp/eminent\_domain.rds}. The
script prints a console summary confirming $n = 312$, $p = 140$, and the
mean and standard deviation of log GDP.

\subsection*{What \texttt{analysis.R} Does}

The analysis script reads \texttt{temp/eminent\_domain.rds} and estimates
Table 2 using four estimators: post-Lasso IV (\texttt{rlassoIV} from
\texttt{hdm}), classical L2-Boosting IV (\texttt{boostSelectZ}, \texttt{post
= FALSE}), post-L2-Boosting IV (\texttt{boostSelectZ}, \texttt{post = TRUE}),
and orthogonal L2-Boosting IV (\texttt{orthoboostSelectZ}), all from the
\texttt{newboost} package. For Table 3, the script runs 500 Monte Carlo
iterations, calling \texttt{DGP()} from \texttt{code/DGP.R} to generate
each draw and applying the same four estimators. Both tables are written
as \LaTeX\ fragments to \texttt{output/tables/table2.tex} and
\texttt{output/tables/table3.tex}. The script prints a side-by-side
console summary comparing replicated values to the paper's targets.

% =============================================================================
\section{Replication Results}
% =============================================================================

%\begin{table}[h]
\centering
\caption{Replication of Table 2: Effect of Federal Appellate Takings Law Decisions
on Economic Outcomes. From Luo and Spindler (2017, AER P\&P).}
\begin{tabular*}{\textwidth}{@{\extracolsep{\fill}}lcccc}
\hline
 & post-Lasso & BA & post-BA & oBA \\
\hline
$\hat{\beta}$ & $0.005$ & $0.005$ & $0.004$ & $0.008$ \\
se             & $0.012$ & $0.007$ & $0.006$ & $0.006$ \\
\hline
\end{tabular*}
\begin{minipage}{\textwidth}
\footnotesize Note: Paper reports $\hat{\beta}$ = 0.005, 0.005, 0.004, 0.008
and se = 0.012, 0.007, 0.006, 0.006.
\end{minipage}
\label{tab:table2}
\end{table}


%\input{../output/tables/table3.tex}

For Table 2, all four estimators produce small, positive point estimates of
the effect of property rights protection on log GDP per capita. The
boosting-based estimators replicate the post-Lasso sign and approximate
magnitude, with BA, post-BA, and oBA returning standard errors modestly
smaller than post-Lasso. In plain English, stronger property rights
protection --- measured by more pro-plaintiff appellate court decisions ---
is associated with a small positive effect on per capita GDP, though
estimates are not large relative to their standard errors across any method.
The economic conclusion is unchanged across all four estimators.

For Table 3, all methods exhibit upward bias under the DGP, with post-Lasso
showing the lowest bias and post-BA the highest. Rejection probabilities are
below the nominal 5\% level for all methods, indicating undersized tests.
The boosting variants perform comparably to post-Lasso in terms of rejection
rates, with BA and oBA performing similarly to each other.

Compared to the original paper, any discrepancies in the replicated estimates
are most likely attributable to differences in the version of the
\texttt{newboost} package, which is unavailable on CRAN and has no documented
version history.

% =============================================================================
\section{What I Learned}
% =============================================================================

The easiest part of the replication was obtaining the data. The
\texttt{EminentDomain} dataset is bundled directly in the \texttt{hdm} package,
so there was no download step, no data cleaning decisions to make, and no
ambiguity about which version of the data was used. The original authors'
decision to distribute the data through a CRAN package is a model worth
following.

The hardest part was the \texttt{newboost} package dependency. The package
is not on CRAN and must be installed from R-Forge or requested from the
authors directly. R-Forge availability is intermittent and no version number
or changelog is documented, making it impossible to verify that the installed
version matches the one used in the original paper. This brings about a concern for future reproduction, as code that depends on a non-CRAN package can
become difficult to run within a few years of publication regardless of how
clean the rest of the replication package is.

Another lesson learned relates to the pipeline structure. The original replication package
contains no Makefile, no directory structure enforcing separation of inputs
and outputs, and no README beyond the brief file description. Reproducing the
paper required manually tracing dependencies across three R scripts
(\texttt{Sim\_AER\_V3.R}, \texttt{DGP.R}, \texttt{helper.R}) with no
automation. Building this replication with the Gentzkow--Shapiro
directory structure and a Makefile made the dependency graph explicit and
revealed variable dependencies in the DGP that could have been easy to break silently. Had a minimal \texttt{run\_all.sh} been provided, the replication effort could
have been substantially reduced.

% =============================================================================
\bibliographystyle{apalike}
\bibliography{references}
% =============================================================================

\end{document}